\section{Control de Formaciones en Sistemas Holonómicos}
\subsection{Error de sincronización de Koulong y Pakniyat}
Punto de partida, tomado de Armel et al.~\cite{koulong2024distributed}.
\begin{equation*}
    e_i =
    \underbrace{-\nu_1 \sum_{j \in \mathcal{N}_i} a_{ij}\big[(\mathbf{x}_i-\boldsymbol{\psi}_i)-(\mathbf{x}_j-\boldsymbol{\psi}_j)\big]}_{\text{\footnotesize Consenso relativo entre vecinos}}
    \;\;
    \underbrace{-\nu_2\, b_i^0\big[(\mathbf{x}_i-\boldsymbol{\psi}_i)-(\mathbf{x}_0-\boldsymbol{\psi}_0)\big]}_{\text{\footnotesize Acoplamiento a la referencia}}
\end{equation*}

{\footnotesize
\noindent
\begin{tabular}{l@{\ }l@{\qquad}l@{\ }l}
$e_i$ & error de sincronización & $\mathcal{N}_i$ & vecinos del agente $i$ \\
$\mathbf{x}_i,\mathbf{x}_j$ & posición de $i$ y $j$ & $a_{ij}$ & peso de comunicación $i$--$j$ \\
$\mathbf{x}_0$ & posición de la referencia & $b_i^0$ & $1$ si $i$ percibe la referencia \\
$\boldsymbol{\psi}_i,\boldsymbol{\psi}_j,\boldsymbol{\psi}_0$ & desplazamiento en la formación & $\nu_1,\nu_2$ & ganancias consenso/acoplamiento \\
\end{tabular}
}
